\begin{pastexam}[label={px:ncu-114-dsa-01}]{中央資工所 114 學年度 第 1 題}
（複選；每題 5 分，全部答對才給分，答錯倒扣 1 分）

An empty 11-bucket hash table has hash function h(k)=k\%11 and adopts quadratic probing (i.e.,
(h(k)+i$^2$)\%11, (h(k)$-$i$^2$)\%11 for i=1, 2, \ldots). Below is the resulting hash table after
inserting 8 items.
\begin{center}
\begin{tabular}{|l|l|}
\hline
0 & \\ \hline
1 & \\ \hline
2 & 25 \\ \hline
3 & 59 \\ \hline
4 & 37 \\ \hline
5 & 27 \\ \hline
6 & 16 \\ \hline
7 & 17 \\ \hline
8 & 15 \\ \hline
9 & 38 \\ \hline
10 & \\ \hline
\end{tabular}
\end{center}
Which of the following is/are the possible input sequence of these 8 items?
\begin{enumerate}[label=(\alph*)]
\item 27 37 16 17 59 15 25 38
\item 27 37 16 59 15 25 38 17
\item 37 27 16 59 25 15 17 38
\item 27 16 37 38 59 15 17 25
\end{enumerate}
\end{pastexam}

\begin{pastsoln}
答案：(a)(b)(c)(d)，四個都可能。

探測順序照題目：$h, h + 1, h - 1, h + 4, h - 4, \ldots$（都取 $\bmod 11$），這是正負交替的二次探測
（資結書第~12.3~節）。 先算每個鍵的 $h(k)$，再看它最後的格子是探測序列的第幾個位置，就知道「它插入
時哪些格子已經有鍵」：
\begin{center}
\small
\begin{tabular}{cccll}
\toprule
鍵 & $h(k)$ & 所在格 & 探測經過 & 插入前必須已經有鍵的格 \\
\midrule
$27$ & $5$ & $5$ & $5$ & 無 \\
$37$ & $4$ & $4$ & $4$ & 無 \\
$16$ & $5$ & $6$ & $5, 6$ & $5$（$27$） \\
$17$ & $6$ & $7$ & $6, 7$ & $6$（$16$） \\
$59$ & $4$ & $3$ & $4, 5, 3$ & $4$（$37$）、$5$（$27$） \\
$15$ & $4$ & $8$ & $4, 5, 3, 8$ & $4, 5, 3$（$37, 27, 59$） \\
$25$ & $3$ & $2$ & $3, 4, 2$ & $3$（$59$）、$4$（$37$） \\
$38$ & $5$ & $9$ & $5, 6, 4, 9$ & $5, 6, 4$（$27, 16, 37$） \\
\bottomrule
\end{tabular}
\end{center}
$15$ 的探測：$4 + 1 = 5$，$4 - 1 = 3$，$4 + 4 = 8$。 $25$：$3 + 1 = 4$，$3 - 1 = 2$。 $38$：$5 + 1 = 6$，
$5 - 1 = 4$，$5 + 4 = 9$。 每個鍵停下的格子前面那些格子必須已經有鍵，而停下的格子本身當時是空的
（每個格子只放一個鍵，後放的鍵不會搶走先放的格子，所以「當時是空的」自動成立）。

把右邊一欄寫成先後條件：
\[
27 \prec 16 \prec 17,\quad 37, 27 \prec 59,\quad 37, 27, 59 \prec 15,\quad 59, 37 \prec 25,\quad
27, 16, 37 \prec 38,
\]
$a \prec b$ 表示 $a$ 要比 $b$ 早插入。 一個輸入順序能產生這張表，若且唯若它滿足這些條件。 逐項檢查：
\begin{itemize}
\item (a) 27 37 16 17 59 15 25 38：$16$ 在 $27$ 後、$17$ 在 $16$ 後；$59$ 在 $37, 27$ 後；$15$ 在
$59$ 後；$25$ 在 $59$ 後；$38$ 在最後。 全部滿足，可能。
\item (b) 27 37 16 59 15 25 38 17：$17$ 放到最後，它只要求 $16$ 先到，仍然滿足；其他條件和 (a) 一樣。
可能。
\item (c) 37 27 16 59 25 15 17 38：$59$ 前面有 $37, 27$；$25$ 前面有 $59, 37$；$15$ 前面有 $37, 27,
59$；$38$ 最後。 可能。
\item (d) 27 16 37 38 59 15 17 25：$38$ 是第 $4$ 個，前面已經有 $27, 16, 37$，探測 $5, 6, 4$ 都有鍵，
落在 $9$；$59$ 前有 $37, 27$；$15$ 前有 $59$；$17$ 前有 $16$；$25$ 前有 $59, 37$。 可能。
\end{itemize}
四個順序照著插入一遍，得到的表都和題目一樣，所以全選。
\end{pastsoln}
